\section{Related Work}

\textbf{Long-context reasoning and verifiable post-training.}
Long-context benchmarks show that nominal context capacity does not guarantee reliable retrieval, integration, or reasoning over evidence distributed across long inputs \citep{bai2024longbench,hsieh2024ruler,liu2024lost,goldman2024really,pmlr-v267-modarressi25a}. Corresponding post-training work improves these capabilities through long-instruction data, length-aware optimization, reinforcement-learning curricula, and task-specific feedback \citep{bai2024longalign,wan2025qwenlong,wang2026loongrl,chen2026longrlvr,zhang2025longreward}. These studies improve the construction of long-context data, curricula, and rewards. Our setting instead begins with an existing task verifier and asks how its response-level assessment should interact with dense token-level OPD guidance.

\textbf{Teacher-centric on-policy distillation.}
MiniLLM and GKD distill on student-generated sequences, exposing the teacher to the student's own generation distribution \citep{gu2024minillm,agarwal2024policy}. ExOPD extrapolates along a teacher--reference direction \citep{yang2026learning}, FiRe-OPD filters trajectories and reweights informative tokens \citep{li2026filter}, and PowerOPD bounds sampled-token rewards for stability \citep{zhao2026poweropd}. Complementary analyses examine teacher--student compatibility and whether large token-level disagreement is learnable, while token-selective distillation concentrates supervision on selected positions \citep{li2026rethinking,wang2026not,kim2026explain}. This line improves how teacher-derived supervision is formed, selected, or optimized. It does not, however, explicitly compare the resulting trajectory-level OPD assessment with a task-specific verifier reward.

\textbf{Verifier-aware on-policy distillation.}
Recent methods connect verifier outcomes to teacher supervision through several interfaces. SCOPE routes correct responses to weighted maximum likelihood and incorrect responses to teacher-perplexity-weighted OPD \citep{zheng2026scope}. MOPD conditions teacher supervision on successful and failed peer responses \citep{yu2026multi}. Uni-OPD calibrates trajectory returns using an outcome-class margin \citep{hou2026uni}. SG-OPD gates token updates by sign consistency and supplements early training with verifier-endorsed teacher rollouts \citep{xu2026sg}. Reward-Weighted OPD selects verifier-passable rollouts and weights their dense teacher supervision \citep{zou2026reward}. These methods introduce verifier feedback through objective routing, teacher conditioning, return calibration, token gating, or trajectory weighting. GC-OPD instead leaves the teacher pass unchanged and forms a response-specific residual from group-relative verifier and OPD assessments. It retains graded within-group differences and the original token-level OPD advantage.

\textbf{Trajectory-level feedback and token-level credit.}
A terminal verifier supplies one response-level reward without specifying its tokenwise influence. Process-supervision methods use human or automatically constructed step labels \citep{lightman2024let,wang2024math}, while VinePPO estimates action advantages from auxiliary continuations \citep{pmlr-v267-kazemnejad25a}. These approaches obtain finer-grained evidence through additional labels or sampling. GC-OPD requires neither step labels nor auxiliary continuations. RACA instead allocates the trajectory-level residual with bounded weights derived from relative OPD advantages without inferring token correctness.
